\begin{abstract}
En este trabajo se presenta un método que evalúa de manera automática las soluciones vecinas de cualquier Problema de Enrutamiento de Vehículos (VRP). La estructura que se propone permite corregir deficiencias del Grafo de Evaluación desarrollado en otra tesis de licenciatura. A esta propuesta se le aplica una cinta computacional basada en la idea que se utiliza en la Diferenciación Automática. Las soluciones vecinas podrán evaluarse sin tener que implementar un código específico para ello y el uso de la cinta garantiza que se evalúen solo las partes de la solución inicial que se modificaron para obtener la nueva solución. Esta propuesta es extensible para cualquier variante de VRP, solo se necesita el código de evaluación de la primera solución.
\end{abstract}

\begin{enabstract}
In this work, a method that automatically evaluates the neighboring solutions of any Vehicle Routing Problem (VRP) is presented. The proposed structure makes it possible to correct deficiencies in the Evaluation Graph developed in another bachelor's thesis. To this proposal a computational tape is applied based on the idea that Automatic Differentiation is used. Neighboring solutions can be evaluated without having to implement specific code for this, and the use of the tape ensures that only those parts of the initial solution that were modified to obtain the new solution are evaluated. This proposal is extensible for any VRP variant, only the evaluation code of the first solution is needed.
\end{enabstract}